%%%%%
%
\section{D�rivation temporelle}

\itemstyle{Op\'erateur~:} \operfmt{Deriv\_Temps}
\index{Deriv\_Temps}

%%%%%OBJET%%%%
\itemstyle{Objet~:}
Calcule la d�riv�e centr�e discr�te, par rapport au temps, d'une variable 
de type \texttt{float} �chantillonn�e r�guli�rement.

%%%%%ARGUMENTS%%%%
\bigskip
%%%%%ENTREES%%%%
\begin{tabarg}
\multicolumn{3}{l}{\titretab {Entr\'ees}}\\
\hline
$X$ & Vecteur, \texttt{float} & Variable (unit� indiff�rente)\\
$f_e$ & Coeff., \texttt{float} & Fr�quence d'�chantillonnage (Hz)\\
\end{tabarg}

\medskip
%%%%%SORTIES%%%%
\begin{tabarg}
\multicolumn{3}{l}{\titretab {Sortie}}\\
\hline
$X_d$ & Vecteur, \texttt{float} & D�riv�e temporelle de $X$ (unit� de $X$ . s$^{-1}$ )\\
\end{tabarg}


%%%%%FORMULE%%%%
\itemstyle{Formule}

%%Deriv_Temps
\[
X_d[i] = \frac{(X[i+1] - X[i-1]) \times f_e }{2}
\]



%%%%%AUTEUR%%%%
\itemstyle{Commentaire~:}
Renvoie une valeur manquante pour les vecteurs de taille 1.

Le calcul de la premi�re et de la derni�re valeur n'est pas centr�.

Si l'�chantillonnage de la variable n'est pas absolument r�gulier, 
utiliser plut�t l'op�rateur \hyperref[OP-Differ_centre_XfYd]{\operfmt{Differ\_centre\_XfYd}}.

%%%%%AUTEUR%%%%
\itemstyle{Auteur~:}
CNRM/GMEI/TRAMM

%%%%%REFERENCE%%%%
\itemstyle{R\'ef\'erences~:}
N�ant
